\documentclass[a4paper,11pt]{article}
\usepackage[utf8]{inputenc}
\usepackage[T1]{fontenc}
\usepackage[french]{babel}
\usepackage{lmodern}
\usepackage{amsmath,amssymb,amsthm}
\usepackage{mathrsfs}
\usepackage{enumitem}
\usepackage{multicol}

\setlist[enumerate]{label=\textbf{\textcolor{Couleur8}{\arabic*.}}, left=1em}

% Si vous voulez aussi modifier les listes enumerate imbriquées :
\setlist[enumerate,2]{label=\textcolor{Couleur8}{\alph*)}}
\setlist[enumerate,3]{label=\textcolor{Couleur8}{\roman*)}}
\setlist[enumerate,4]{label=\textcolor{Couleur8}{\Alph*)}}
\usepackage{graphicx}
\usepackage{xcolor}
\usepackage{tcolorbox}
\usepackage{geometry}
\usepackage{fancyhdr}
\usepackage{titlesec}
\usepackage{cancel}
\usepackage{ifthen}
\usepackage{tikz}
\usetikzlibrary{arrows.meta}
\usepackage{tkz-tab}
\usepackage{eurosym}

% OPTION PROF/ÉLÈVE
% Décommentez UNE des deux lignes suivantes :
\newboolean{prof}\setboolean{prof}{true}   % Version professeur (avec corrections des devoirs)
%\newboolean{prof}\setboolean{prof}{false}  % Version élève (sans corrections des devoirs)

\definecolor{Couleur1}{HTML}{4A5D7A} %Th
\definecolor{Couleur2}{HTML}{A5B5D6} %Prop
\definecolor{Couleur3}{HTML}{A8C68A} %Méthode
\definecolor{Couleur6}{HTML}{8A9CC1} %Def
\definecolor{Couleur7}{HTML}{E6B459} %Rq Voc
\definecolor{Couleur8}{HTML}{D36740} %Titres
\definecolor{correction}{HTML}{D36740} % Couleur pour les corrections
\definecolor{coopmathscorr}{HTML}{F15929} % Couleur corrections coopmaths

% Configuration des marges
\geometry{
	top=2cm,
	bottom=2cm,
	inner=1.5cm,
	outer=1.5cm,
}

% Police
\renewcommand{\familydefault}{\sfdefault}

% Compteurs
\newcounter{seancenum}
\newcounter{seancenumD}
\setcounter{seancenum}{1}
\setcounter{seancenumD}{1}

% Configuration des en-têtes et pieds de page
\pagestyle{fancy}
\fancyhf{}
\ifthenelse{\boolean{prof}}{
    \fancyhead[L]{\textcolor{Couleur8}{\textbf{Corrections Automatismes - Classe de seconde - Version Professeur}}}
}{
    \fancyhead[L]{\textcolor{Couleur8}{\textbf{Corrections Automatismes - Séance 7 - Classe de seconde}}}
}
\fancyhead[R]{\textcolor{Couleur8}{\textbf{2025-2026}}}
\fancyfoot[L]{\textcolor{Couleur8}{Mme Bertail et Mme Pinto}}
\fancyfoot[R]{\textcolor{Couleur8}{\thepage}}
\renewcommand{\headrulewidth}{1pt}
\renewcommand{\headrule}{\hbox to\headwidth{\color{Couleur8}\leaders\hrule height \headrulewidth\hfill}}

% Environnements personnalisés
\newtcolorbox{seance}[1]{
    colback=Couleur6!10,
    colframe=Couleur1!65,
    fonttitle=\bfseries\large,
    title=Correction Séance \theseancenum #1,
    left=5mm,
    right=5mm,
    top=3mm,
    bottom=3mm,
    arc=0mm,
    after=\stepcounter{seancenum}
}

\newtcolorbox{seanceD}[1]{
    colback=Couleur6!5,
    colframe=Couleur6!45,
    fonttitle=\bfseries\large,
    title=Correction Devoirs - Séance \theseancenumD #1,
    left=5mm,
    right=5mm,
    top=3mm,
    bottom=3mm,
    arc=0mm,
    after=\stepcounter{seancenumD}
}

\begin{document}

\setcounter{seancenum}{7}
\newpage
\begin{seance}{} %7

\begin{enumerate}
\item Pour passer de $20\,\%$ à $100\,\%$, on multiplie par $5$. \\
$\begin{aligned}
20\,\% \text{ de } 150 &= 30\\
100\,\% \text{ de } 150 &= 5 \times 30\\
N &= 150
\end{aligned}$

Donc $p = {\color{correction}\boldsymbol{20}}$.

\item Pour calculer $25\,\%$ de $240$, on effectue le calcul $0{,}25 \times 240$ soit encore ${\color{correction}\boldsymbol{\dfrac{25 \times 240}{100}}}$.

La bonne réponse est la réponse {\bfseries \color{correction}B}.

\item $15\,\%$ de $400 = 0{,}15 \times 400 = {\color{correction}\boldsymbol{60}}$

Donc ${\color{correction}\boldsymbol{60}}$ élèves étudient l'Italien.

\item Pour passer de $30\,\%$ à $100\,\%$, on divise par $30$ puis on multiplie par $100$ (ou on multiplie par $\dfrac{100}{30} = \dfrac{10}{3}$).
$\begin{aligned}
30\,\% \text{ de } N &= 45\\
100\,\% \text{ de } N &= \dfrac{100}{30} \times 45\\
N &= \dfrac{10}{3} \times 45 = {\color{correction}\boldsymbol{150}}
\end{aligned}$

\item Augmenter de $8\,\%$ revient à multiplier par $1 + \dfrac{8}{100}$.
Ainsi, le coefficient multiplicateur associé à une augmentation de $8\,\%$ est $1 + 0{,}08$, soit ${\color{correction}\boldsymbol{1{,}08}}$.

\item L'augmentation est de $21$ euros sur un total de $60$ euros.
Le pourcentage d'augmentation est donné par le quotient : $\dfrac{21}{60} = \dfrac{7}{20} = 0{,}35 = {\color{correction}\boldsymbol{35}}\,\%$.

\end{enumerate}

\end{seance}

\end{document}